% =====================================================================
%  整理者　macanine <macanine@163.com>
%  仓库　　github.com/macanine/zju-linear-algebra-papers
%  来源　　sources/2024-2025-秋冬-期中-甲.pdf 全册（参考答案另见 sources/2024-2025-秋冬-期中-甲-参考答案.pdf）
% =====================================================================

\begin{problem}{}{8.5cm}
设矩阵 $A$ 的伴随矩阵
$A^{*}=\begin{pmatrix}1&0&0&0\\0&1&0&0\\1&0&1&0\\0&-3&0&8\end{pmatrix}$，
且 $ABA^{-1}=BA^{-1}+3E$，其中 $E$ 为单位矩阵，求 $B$。
\end{problem}

\begin{problem}{}{8.5cm}
设向量组 $\alpha_1=(2,-1,3,4)^{T}$，$\alpha_2=(1,1,1,1)^{T}$，$\alpha_3=(1,2,1,1)^{T}$，$\alpha_4=(5,3,6,7)^{T}$。
求此向量组的秩和一个极大线性无关组，并将其余向量用此极大线性无关组线性表示。
\end{problem}

\begin{problem}{}{8.5cm}
设四元非齐次线性方程组 $AX=b$ 的系数矩阵的秩为 $3$，已知 $\eta_1,\eta_2,\eta_3$ 是它的三个解向量，且
$\eta_1=(2,3,4,5)^{T}$，$\eta_2+\eta_3=(1,2,3,4)^{T}$，求该线性方程组的通解。
\end{problem}

\begin{problem}{}{10cm}
在欧氏空间中，设有两组基
\begin{align*}
(\mathrm{I}):\quad & \varepsilon_1=(1,1,1)^{T},\ \varepsilon_2=(1,1,0)^{T},\ \varepsilon_3=(1,0,0)^{T};\\
(\mathrm{II}):\quad & \eta_1=(-5,2,1)^{T},\ \eta_2=(4,-3,4)^{T},\ \eta_3=(3,-2,2)^{T}.
\end{align*}
（1）求基（I）到基（II）的过渡矩阵；\quad（2）用施密特正交化方法将基（I）改造成标准正交基。
\end{problem}

\begin{problem}{}{8.5cm}
（1）试证明：$n$ 阶反对称矩阵可逆的必要条件是 $n$ 为偶数；\quad（2）该条件是充分条件吗？为什么？
\end{problem}

\begin{problem}{}{8cm}
解矩阵方程 $AX=A+2X$，其中
$A=\begin{pmatrix}4&2&3\\1&1&0\\-1&2&3\end{pmatrix}$。
\end{problem}

\begin{problem}{}{8.5cm}
求下列向量组的秩和一个极大线性无关组
\begin{align*}
\alpha_1&=(6,4,1,-1,2), & \alpha_2&=(1,0,2,3,-4),\\
\alpha_3&=(1,4,-9,-16,22), & \alpha_4&=(7,1,0,-1,3).
\end{align*}
\end{problem}

\begin{problem}{}{10cm}
在线性空间 $P^{3}$ 中，设有两组基
\begin{align*}
(\mathrm{I}):\quad & \varepsilon_1=(1,2,1)^{T},\ \varepsilon_2=(2,3,3)^{T},\ \varepsilon_3=(3,7,1)^{T};\\
(\mathrm{II}):\quad & \beta_1=(3,1,4)^{T},\ \beta_2=(5,2,1)^{T},\ \beta_3=(1,1,-6)^{T}.
\end{align*}
（1）求基（I）到基（II）的过渡矩阵；
（2）若向量 $\alpha$ 在基（I）下的坐标为 $X=(0,1,-1)^{T}$，求 $\alpha$ 在基（II）下的坐标。
\end{problem}

\begin{problem}{}{8cm}
设 $\alpha=(a_1,a_2)$，$\beta=(b_1,b_2)$ 为二阶实空间 $R^{2}$ 中的任意两个向量，问：$R^{2}$ 对以下所规定的内积是否构成欧氏空间？
\begin{enumerate}[label=(\arabic*),leftmargin=2.2em,itemsep=0.25em,topsep=0.3em]
  \item $(\alpha,\beta)=a_1b_2+a_2b_1$；
  \item $(\alpha,\beta)=(a_1+a_2)b_1+(a_1+2a_2)b_2$。
\end{enumerate}
\end{problem}

\begin{problem}{}{8.5cm}
设
$\beta_1=\alpha_2+\alpha_3+\cdots+\alpha_m$，
$\beta_2=\alpha_1+\alpha_3+\cdots+\alpha_m$，
$\cdots$，
$\beta_m=\alpha_1+\alpha_2+\cdots+\alpha_{m-1}$，
其中 $m>1$。证明：向量组 $\beta_1,\beta_2,\cdots,\beta_m$ 与 $\alpha_1,\alpha_2,\cdots,\alpha_m$ 有相同的秩。
\end{problem}

\begin{problem}{}{8cm}
设 $A,B$ 均为 $n$ 阶可逆方阵，证明：如果 $A+B$ 可逆，则 $A^{-1}+B^{-1}$ 也可逆，并求其逆方阵。
\end{problem}

\begin{problem}{}{9cm}
解线性方程组
\[
\begin{cases}
x_1+x_2-2x_3+3x_4=0,\\
2x_1+x_2-6x_3+4x_4=-1,\\
3x_1+2x_2-8x_3+7x_4=-1,\\
x_1-x_2-6x_3-x_4=-2.
\end{cases}
\]
\end{problem}

\begin{problem}{}{7cm}
计算下列行列式
\[
\begin{vmatrix}
1+x&1&1&1\\
1&1-x&1&1\\
1&1&1+y&1\\
1&1&1&1-y
\end{vmatrix}.
\]
\end{problem}

\begin{problem}{}{8cm}
计算下列行列式
\[
\begin{vmatrix}
1&2&3&\cdots&n\\
2&3&4&\cdots&1\\
3&4&5&\cdots&2\\
\vdots&\vdots&\vdots&&\vdots\\
n&1&2&\cdots&n-1
\end{vmatrix}.
\]
\end{problem}

\begin{problem}{}{9cm}
设有线性方程组
\[
\begin{cases}
x_1+a_1x_2+a_1^{2}x_3=a_1^{3},\\
x_1+a_2x_2+a_2^{2}x_3=a_2^{3},\\
x_1+a_3x_2+a_3^{2}x_3=a_3^{3},\\
x_1+a_4x_2+a_4^{2}x_3=a_4^{3}.
\end{cases}
\]
（1）证明：当 $a_1,a_2,a_3,a_4$ 互异时，此方程组无解；
（2）设 $a_1=a_3=k$，$a_2=a_4=-k$，且 $\eta=(-1,1,1)^{T}$ 为此方程组的解，求此方程组的全部解。
\end{problem}

\begin{problem}{}{9.5cm}
试证明
\[
D_n=\begin{vmatrix}
1&2&3&\cdots&n-1&n\\
x&1&2&\cdots&n-2&n-1\\
x&x&1&\cdots&n-3&n-2\\
\vdots&\vdots&\vdots&&\vdots&\vdots\\
x&x&x&\cdots&1&2\\
x&x&x&\cdots&x&1
\end{vmatrix}=(-1)^{n}\bigl[(x-1)^{n}-x^{n}\bigr].
\]
\end{problem}

\begin{problem}{}{6.5cm}
已知
$\begin{vmatrix}a_{11}&a_{12}&a_{13}\\a_{21}&a_{22}&a_{23}\\a_{31}&a_{32}&a_{33}\end{vmatrix}=1$，
计算行列式
$\begin{vmatrix}3a_{11}&4a_{11}-a_{12}&-a_{13}\\3a_{21}&4a_{21}-a_{22}&-a_{23}\\3a_{31}&4a_{31}-a_{32}&-a_{33}\end{vmatrix}$。
\end{problem}

\begin{problem}{}{10cm}
问 $a,b$ 取何值时，线性方程组
\[
\begin{cases}
x_1+x_2+x_3+x_4=0,\\
x_2+2x_3+2x_4=1,\\
-x_2+(a-3)x_3-2x_4=b,\\
3x_1+2x_2+x_3+ax_4=-1
\end{cases}
\]
无解，有唯一解或有无穷解？在有解时，求出解。
\end{problem}

\begin{problem}{}{8.5cm}
解线性方程组
\[
\begin{cases}
(1+a_1)x_1+x_2+x_3+\cdots+x_n=b_1,\\
x_1+(1+a_2)x_2+x_3+\cdots+x_n=b_2,\\
\qquad\cdots\qquad\cdots\qquad\cdots\\
x_1+x_2+x_3+\cdots+(1+a_n)x_n=b_n,
\end{cases}
\]
其中 $a_i\ne 0$，$i=1,2,\cdots,n$，且 $\dfrac{1}{a_1}+\dfrac{1}{a_2}+\cdots+\dfrac{1}{a_n}\ne -1$。
\end{problem}

\begin{problem}{}{9cm}
计算 $n+1$ 阶行列式
\[
D_{n+1}=\begin{vmatrix}
0&0&0&\cdots&0&1&-1\\
0&0&0&\cdots&1&-1&b_1\\
0&0&0&\cdots&-1&0&b_2\\
\vdots&\vdots&\vdots&&\vdots&\vdots&\vdots\\
0&1&-1&\cdots&0&0&b_{n-2}\\
1&-1&0&\cdots&0&0&b_{n-1}\\
-1&0&0&\cdots&0&0&b_n
\end{vmatrix}.
\]
\end{problem}

\begin{problem}{}{8cm}
证明：若某行列式的某行元素全为 $1$，则这个行列式全部代数余子式的和等于该行列式的值。
\end{problem}
